% Part: proof-theory
% Chapter: cut-elimination
% Section: ce-topmost

\documentclass[../../../include/open-logic-section]{subfiles}

\begin{document}

\olfileid[fr]{pt}{cut}{top}

\olsection{Élimination des coupures les plus hautes}

On élimine d'abord une coupure dont les deux prémisses ont
des preuves sans coupure, puis on répète cette opération.
















\begin{defn}
Considérons une !!{proof}~$\pi$ qui se termine par une règle
\CutCS
\[
\AxiomC{}
\RightLabel{$\pi_1$}
\Deduce$\Gamma \fCenter \Delta, !A$
\AxiomC{}
\RightLabel{$\pi_2$}
\Deduce$!A, \Gamma \fCenter \Delta$
\RightLabel{\CutCS}
\BinaryInf$\Gamma \fCenter \Delta$
\DisplayProof
\]
et ne contient aucune autre inférence \CutCS.
La \emph{hauteur de coupure} de~$\pi$ est
$\cheight{\pi}=\pheight{\pi_1}+\pheight{\pi_2}$.
Son \emph{rang de coupure} est
$\cutrank{\pi}=\depth{!A}$.
\end{defn}

\begin{lem}\ollabel{lem:cut-adm-G3c}
La règle \CutCS est admissible dans~\Log{G3c}.
Autrement dit, si $\pi$ est une !!{proof} qui se termine
par une unique inférence~\CutCS et ne comporte ailleurs
que des règles de~$\Log{G3c}$, il existe une !!{proof}~$\pi'$
du même séquent final dans~\Log{G3c}.
\end{lem}

\begin{proof}
Nous raisonnons par double récurrence lexicographique,
d'abord sur le rang de coupure, puis sur sa !!{height}.
La !!{proof}~$\pi$ se termine par
\[
\AxiomC{}
\RightLabel{$\pi_1$}
\Deduce$\Gamma \fCenter \Delta, !A$
\AxiomC{}
\RightLabel{$\pi_2$}
\Deduce$!A, \Gamma \fCenter \Delta$
\RightLabel{\CutCS}
\BinaryInf$\Gamma \fCenter \Delta$
\DisplayProof
\]

Au cas initial, $\cheight{\pi}=0$ et $\cutrank{\pi}=0$.
Puisque $\cheight{\pi}=\pheight{\pi_1}+\pheight{\pi_2}=0$,
les séquents $\Gamma\Sequent\Delta,!A$ et
$!A,\Gamma\Sequent\Delta$ sont tous deux des axiomes.
Pour les axiomes d'identité, deux cas se présentent :
(a)~$!A$ n'est pas principal dans
l'un d'eux, ou (b)~$!A$ est principal dans les deux.
Nous montrerons ci-dessous que, dans chaque cas,
$\Gamma\Sequent\Delta$ est un axiome, sans utiliser
l'hypothèse de récurrence. Les axiomes de fausseté sont
traités dans la variante qui suit.

Distinguons maintenant les cas suivant que $\pi_1$ et $\pi_2$
sont des axiomes ou se terminent par une inférence, et suivant
que $!A$ y est principal ou non.
Les cas A et~B couvrent le cas initial.
À l'étape de récurrence, $\cheight{\pi}>0$ ou
$\cutrank{\pi}>0$, ou les deux. L'hypothèse de récurrence
permet de supprimer l'unique coupure finale d'une !!{proof}
si son rang est $=\cutrank{\pi}$ et sa hauteur de coupure
est $<\cheight{\pi}$, ou si son rang est~$<\cutrank{\pi}$.
Dans cette dernière situation, sa hauteur de coupure peut
être quelconque. Le cas $\cheight{\pi}=0$ est couvert
par A et~B, et les cas $\cheight{\pi}>0$ par C--E.

\paragraph{Cas A.}
Une formule $!B$ figure à la fois dans $\Gamma$ et dans~$\Delta$.
Alors $\Gamma\Sequent\Delta$ est un axiome si $!B$ est atomique;
sinon, \olref[seq][ex]{prop:G3c-axioms} en fournit une
!!{proof}~$\pi'$ dans \Log{G3c}, sans \CutCS.
Cela couvre le cas~(a) initial pour les axiomes d'identité :
si les deux prémisses sont des axiomes d'identité et que $!A$ n'est pas principal dans
$\Gamma\Sequent\Delta,!A$ ou dans
$!A,\Gamma\Sequent\Delta$, un autre atome~$!B$ figure
des deux côtés. Cela couvre aussi l'étape de récurrence
lorsqu'une prémisse est un axiome d'identité où $!A$ n'est pas principal,
même si l'autre prémisse est la conclusion d'une règle.

\paragraph{Cas B.}
Les deux prémisses sont des axiomes d'identité et $!A$ est principal,
donc atomique. Dans la prémisse gauche,
$\Gamma\Sequent\Delta,!A$, la formule $!A$ doit aussi
figurer dans $\Gamma$, sans quoi ce séquent ne serait pas
un axiome. De même, la prémisse droite
$!A,\Gamma\Sequent\Delta$ exige que $!A$ figure dans $\Delta$.
La formule atomique $!A$ se trouve donc dans $\Gamma$
et dans $\Delta$; $\Gamma\Sequent\Delta$ est un axiome.
Cela couvre le cas~(b) initial.

\begin{center}
Variante des cas A et B
\end{center}

\paragraph{Cas A.}
Une prémisse de la coupure est un axiome
$!C,\Gamma'\Sequent\Delta',!C$, avec $!C$ atomique.
Si $!C$ n'est pas la !!{formula} coupée~$!A$, il figure
dans $\Gamma$ et dans~$\Delta$.
Le séquent $\Gamma\Sequent\Delta$ est alors un axiome
et constitue la preuve~$\pi'$ recherchée.
Si $!C\ident!A$, deux possibilités se présentent.
Lorsque la prémisse gauche est l'axiome, on a
$!A\in\Gamma$ et $\Gamma=\Gamma',!A$.
La preuve $\pi_2$ porte alors sur
$!A,!A,\Gamma'\Sequent\Delta$;
l'admissibilité de \LeftR{\Contraction},
\olref[seq][inv]{prop:G3c-cont-adm}, donne une preuve
de $\Gamma\Sequent\Delta$.
Lorsque la prémisse droite est l'axiome, on a
$!A\in\Delta$ et $\Delta=\Delta',!A$.
C'est alors $\pi_1$ qui !!{prove}s
$\Gamma\Sequent\Delta',!A,!A$.
L'admissibilité de $\RightR{\Contraction}$ donne
$\Gamma\Sequent\Delta$.

\paragraph{Cas B.}
Une prémisse est un axiome de la forme
$\lfalse,\Gamma'\Sequent\Delta'$.
Si c'est la prémisse gauche, $\lfalse\in\Gamma$
et $\Gamma\Sequent\Delta$ est aussi un axiome.
Si c'est la prémisse droite, soit $\lfalse\in\Gamma$,
et le résultat est encore immédiat, soit $!A\ident\lfalse$.
Dans ce dernier cas, $\pi_1$ !!{prove}s
$\Gamma\Sequent\Delta,\lfalse$.
On obtient une !!{proof} de~$\Gamma\Sequent\Delta$
en supprimant une copie suivie de~$\lfalse$ à droite
dans chaque séquent de~$\pi_1$.
Vérifier ce fait par récurrence sur~$\pheight{\pi_1}$ :
aucune de ces occurrences n'est principale et les axiomes
restent des axiomes.

Supposons désormais que ni A ni B ne s'applique.
On a $\cheight{\pi}>0$ : au moins une prémisse
est la conclusion d'une inférence.

\paragraph{Cas C et D : permutation des coupures.}
Ces cas suivent un même schéma.
Une règle~$R$ précède la coupure \CutCS, et la !!{formula}
coupée~$!A$ n'est pas principale dans cette règle;
une autre formule~$!D$ l'est.
Schématiquement :
\[
\AxiomC{}
\RightLabel{$\pi_1$}
\Deduce$\Gamma \fCenter \Delta', !D, !A$
\AxiomC{}
\RightLabel{$\pi_3$}
\Deduce$!A, \Gamma, \Pi \fCenter \Delta', \Lambda$
\RightLabel{$R$}
\UnaryInf$!A, \Gamma \fCenter \Delta', !D$
\RightSubproofLabel{$\pi_2$}
\RightLabel{\CutCS}
\BinaryInf$\Gamma \fCenter \Delta', !D$
\DisplayProof
\]
Ce diagramme représente seulement une règle droite~$R$
à une prémisse. La formule $!A$ n'est pas principale,
et la !!{formula} principale~$!D$ figure dans le succédent
de la prémisse droite de \CutCS.
Si $!D$ figurait dans l'antécédent, ou si $R$ s'appliquait
dans la prémisse gauche de la coupure, le diagramme serait
différent. Les !!{formula}s de $\Pi$ et de~$\Lambda$
sont les formules actives de~$R$.

Inversons l'ordre de ces deux inférences :
\[
\AxiomC{}
\RightLabel{$\pi_1'$}
\Deduce$\Gamma, \Pi \fCenter \Delta', !D, \Lambda, !A$
\AxiomC{}
\RightLabel{$\pi_3'$}
\Deduce$!A, \Gamma, \Pi \fCenter \Delta', !D, \Lambda$
\RightLabel{\CutCS}
\BinaryInf$\Gamma, \Pi \fCenter \Delta', !D, \Lambda$
\RightLabel{$R$}
\UnaryInf$\Gamma \fCenter \Delta', !D, !D$
\DisplayProof
\]
Les preuves $ \pi_1'$ et $\pi_3'$ sont obtenues à partir
de $\pi_1$ et de $\pi_3$ par affaiblissement sans
augmentation de hauteur.

On dit qu'on \emph{permute la coupure vers le haut à travers~$R$}.
La sous-!!{proof} conduisant à la prémisse droite de la
nouvelle coupure perd une inférence.
Les hauteurs des deux sous-!!{proof}s de cette coupure
ne dépassent donc pas celles de $\pi_1$ et de $\pi_3$;
la hauteur de coupure diminue.
L'hypothèse de récurrence permet de supprimer la coupure.
On réapplique ensuite la règle et on supprime la copie
supplémentaire de~$!D$ par admissibilité de
\RightR{\Contraction}.

\paragraph{Cas C.}
Au moins une des preuves $\pi_1$, $\pi_2$ se termine
par une règle~$R$ à une prémisse où $!A$ n'est pas principal.
Il y a 18 cas : deux choix du côté de la coupure et neuf
choix de règle, à savoir \LeftR{\lnot}, \RightR{\lnot},
\RightR{\lor}, \LeftR{\land}, \RightR{\lif} et les quatre
règles quantifiées. Nous traitons deux exemples où
$!A$ n'est pas principal dans la dernière inférence
de~$\pi_2$ : \RightR{\lif} et \LeftR{\lexists}.
Les autres cas sont analogues et laissés en exercice.

Supposons que $\pi$ se termine par
\[
\AxiomC{}
\RightLabel{$\pi_1$}
\Deduce$\Gamma \fCenter \Delta', !B \lif !C, !A$
\AxiomC{}
\RightLabel{$\pi_3$}
\Deduce$!A, !B, \Gamma \fCenter \Delta', !C$
\RightLabel{$\RightR{\lif}$}
\UnaryInf$!A, \Gamma \fCenter \Delta', !B \lif !C$
\RightSubproofLabel{$\pi_2$}
\RightLabel{\CutCS}
\BinaryInf$\Gamma \fCenter \Delta', !B \lif !C$
\DisplayProof
\]
où $\Delta=\Delta',!B\lif!C$.
L'affaiblissement sans augmentation de hauteur,
\olref[seq][adm]{prop:weak-G3c-adm}, transforme $\pi_1$
en une !!{proof} $\pi_1'$ de
$!B,\Gamma\fCenter\Delta',!B\lif!C,!C,!A$,
en ajoutant $!B$ à gauche et $!C$ à droite.
De même, \olref[seq][adm]{prop:weak-G3c-adm}
transforme~$\pi_3$ en une !!{proof}~$\pi_3'$ de
$!A,!B,\Gamma\fCenter\Delta',!B\lif!C,!C$,
en ajoutant $!B\lif!C$ à droite.
L'hypothèse de récurrence s'applique à
\[
\AxiomC{}
\RightLabel{$\pi_1'$}
\Deduce$!B, \Gamma \fCenter \Delta', !B \lif !C, !C, !A$
\AxiomC{}
\RightLabel{$\pi_3'$}
\Deduce$!A, !B, \Gamma \fCenter \Delta', !B \lif !C, !C$
\RightLabel{\CutCS}
\BinaryInf$!B, \Gamma \fCenter \Delta', !B \lif !C, !C$
\DisplayProof
\]
La !!{formula} coupée est $!A$; le rang est inchangé,
mais la hauteur de coupure a diminué.
On obtient une !!{proof}~$\pi_4$ dans \Log{G3c}
de~$!B,\Gamma\fCenter\Delta',!B\lif!C,!C$.
La règle $\RightR{\lif}$ donne
$\Gamma\fCenter\Delta',!B\lif!C,!B\lif!C$ :
\[
\AxiomC{}
\RightLabel{$\pi_4$}
\Deduce$!B, \Gamma \fCenter \Delta', !B \lif !C, !C$
\RightLabel{$\RightR{\lif}$}
\UnaryInf$\Gamma \fCenter \Delta', !B \lif !C, !B \lif !C$
\DisplayProof
\]
Enfin, \olref[seq][inv]{prop:G3c-cont-adm} supprime la
copie supplémentaire et fournit une !!{proof}~$\pi'$
de $\Gamma\Sequent\Delta$, c'est-à-dire
$\Gamma\fCenter\Delta',!B\lif!C$.

Supposons maintenant que $\pi$ se termine par
\[
\AxiomC{}
\RightLabel{$\pi_1$}
\Deduce$\lexists[x][!B(x)], \Gamma' \fCenter \Delta, !A$
\AxiomC{}
\RightLabel{$\pi_3$}
\Deduce$!A, !B(c), \Gamma' \fCenter \Delta$
\RightLabel{$\LeftR{\lexists}$}
\UnaryInf$!A, \lexists[x][!B(x)], \Gamma' \fCenter \Delta$
\RightSubproofLabel{$\pi_2$}
\RightLabel{\CutCS}
\BinaryInf$\lexists[x][!B(x)], \Gamma' \fCenter \Delta$
\DisplayProof
\]
L'hypothèse de récurrence s'applique de nouveau à la
coupure du diagramme suivant :
\[
\AxiomC{}
\RightLabel{$\pi_1[!B(c) \Sequent {}]$}
\Deduce$\lexists[x][!B(x)], !B(c), \Gamma' \fCenter \Delta, !A$
\AxiomC{}
\LeftLabel{$\pi_3[\lexists[x][!B(x)] \Sequent {}]$}
\Deduce$!A, \lexists[x][!B(x)], !B(c), \Gamma' \fCenter \Delta$
\RightLabel{\CutCS}
\BinaryInf$\lexists[x][!B(x)], !B(c), \Gamma' \fCenter \Delta$
\RightLabel{\LeftR{\lexists}}
\UnaryInf$\lexists[x][!B(x)], \lexists[x][!B(x)], \Gamma' \fCenter \Delta$
\DisplayProof
\]
La condition de variable propre de
\LeftR{\lexists} dans~$\pi_2$ garantit que $c$
n'apparaît ni dans $\lexists[x][!B(x)]$, ni dans
$\Gamma'$, ni dans~$\Delta$.
La même condition est donc satisfaite dans la règle
reconstruite. On termine par la
\olref[seq][inv]{prop:G3c-cont-adm}.

\paragraph{Cas D.}
Au moins une des preuves $\pi_1$, $\pi_2$ se termine
par une règle~$R$ à deux prémisses où $!A$ n'est pas principal.
Il y a six cas. Traitons celui où $\pi_2$ se termine
par~$\RightR{\land}$; les autres sont analogues.

Supposons que $\pi$ se termine par
\[
\AxiomC{}
\RightLabel{$\pi_1$}
\Deduce$\Gamma \fCenter \Delta', !B \land !C, !A$
\AxiomC{}
\RightLabel{$\pi_3$}
\Deduce$!A, \Gamma \fCenter \Delta', !B$
\AxiomC{}
\RightLabel{$\pi_4$}
\Deduce$!A, \Gamma \fCenter \Delta', !C$
\RightLabel{$\RightR{\land}$}
\BinaryInf$!A, \Gamma \fCenter \Delta', !B \land !C$
\RightSubproofLabel{$\pi_2$}
\RightLabel{\CutCS}
\BinaryInf$\Gamma \fCenter \Delta$
\DisplayProof
\]
L'hypothèse de récurrence s'applique aux !!{proof}s
\[
\AxiomC{}
\RightLabel{$\pi_1'$}
\Deduce$\Gamma \fCenter \Delta', !B \land !C, !B, !A$
\AxiomC{}
\RightLabel{$\pi_3'$}
\Deduce$!A, \Gamma \fCenter \Delta', !B \land !C, !B$
\RightLabel{\CutCS}
\BinaryInf$\Gamma \fCenter \Delta', !B \land !C, !B$
\DisplayProof
\]
et
\[
\AxiomC{}
\RightLabel{$\pi_1''$}
\Deduce$\Gamma \fCenter \Delta', !B \land !C, !C, !A$
\AxiomC{}
\RightLabel{$\pi_4'$}
\Deduce$!A, \Gamma \fCenter \Delta', !B \land !C, !C$
\RightLabel{\CutCS}
\BinaryInf$\Gamma \fCenter \Delta', !B \land !C, !C$
\DisplayProof
\]
Les !!{proof}s de \Log{G3c} $\pi_1'$ et $\pi_1''$
proviennent de~$\pi_1$ par affaiblissement sans
augmentation de hauteur,
\olref[seq][adm]{prop:weak-G3c-adm} :
on ajoute respectivement $!B$ et $!C$ à droite.
Les !!{proof}s $\pi_3'$ et~$\pi_4'$ proviennent de
$\pi_3$ et de~$\pi_4$ par le même affaiblissement,
en ajoutant $!B\land!C$ à droite.

L'hypothèse de récurrence fournit des !!{proof}s
$\pi_5$ et~$\pi_6$ dans \Log{G3c} de
$\Gamma\fCenter\Delta',!B\land!C,!B$ et
$\Gamma\fCenter\Delta',!B\land!C,!C$.
On les réunit par la règle $\RightR{\land}$ pour
obtenir $\Gamma\Sequent\Delta',!B\land!C,!B\land!C$ :
\[
\AxiomC{}
\RightLabel{$\pi_5$}
\Deduce$\Gamma \fCenter \Delta', !B \land !C, !B$
\AxiomC{}
\RightLabel{$\pi_6$}
\Deduce$\Gamma \fCenter \Delta', !B \land !C, !C$
\RightLabel{\RightR{\land}}
\BinaryInf$\Gamma \fCenter \Delta', !B \land !C, !B \land !C$
\DisplayProof
\]
On applique enfin \olref[seq][inv]{prop:G3c-cont-adm}
pour obtenir une !!{proof} de $\Gamma\Sequent\Delta$,
c'est-à-dire $\Gamma\fCenter\Delta',!B\land!C$.

\paragraph{Cas E : réduction des coupures.}
Les preuves $\pi_1$ et $\pi_2$ se terminent toutes deux
par des inférences où $!A$ est principal.
Il y a six cas, un pour chaque !!{main operator} possible.
Dans chacun d'eux, on remplace la coupure sur~$!A$
par une ou plusieurs coupures sur ses sous-!!{formula}s,
qui sont les formules actives des inférences précédentes.
On parle de \emph{réduction} ou de \emph{conversion}
de l'inférence de coupure.

\begin{enumerate}
\item Si $!A\ident\lnot!B$, la preuve $\pi$ se termine par
\[
\AxiomC{}
\RightLabel{$\pi_1'$}
\Deduce$!B, \Gamma \fCenter \Delta$
\RightLabel{\RightR{\lnot}}
\UnaryInf$\Gamma \fCenter \Delta, \lnot !B$
\LeftSubproofLabel{$\pi_1$}
\AxiomC{}
\RightLabel{$\pi_2'$}
\Deduce$\Gamma \fCenter \Delta, !B$
\RightLabel{\LeftR{\lnot}}
\UnaryInf$\lnot !B, \Gamma \fCenter \Delta$
\RightSubproofLabel{$\pi_2$}
\RightLabel{\CutCS}
\BinaryInf$\Gamma \fCenter \Delta$
\DisplayProof
\]
L'hypothèse de récurrence permet de supprimer \CutCS de
\[
\AxiomC{}
\RightLabel{$\pi_2'$}
\Deduce$\Gamma \fCenter \Delta, !B$
\AxiomC{}
\RightLabel{$\pi_1'$}
\Deduce$!B, \Gamma \fCenter \Delta$
\RightLabel{\CutCS}
\BinaryInf$\Gamma \fCenter \Delta$
\DisplayProof
\]

\item Si $!A\ident(!B\lor!C)$, la preuve $\pi$ se termine par
\[
\AxiomC{}
\RightLabel{$\pi_1'$}
\Deduce$\Gamma \fCenter \Delta, !B, !C$
\RightLabel{\RightR{\lor}}
\UnaryInf$\Gamma \fCenter \Delta, !B \lor !C$
\LeftSubproofLabel{$\pi_1$}
\AxiomC{}
\RightLabel{$\pi_2'$}
\Deduce$!B, \Gamma \fCenter \Delta$
\AxiomC{}
\RightLabel{$\pi_2''$}
\Deduce$!C, \Gamma \fCenter \Delta$
\RightLabel{\LeftR{\lor}}
\BinaryInf$!B \lor !C, \Gamma \fCenter \Delta$
\RightSubproofLabel{$\pi_2$}
\RightLabel{\CutCS}
\BinaryInf$\Gamma \fCenter \Delta$
\DisplayProof
\]
Considérons la !!{proof} qui se termine par \CutCS :
\[
\AxiomC{}
\RightLabel{$\pi_1'$}
\Deduce$\Gamma \fCenter \Delta, !B, !C$
\AxiomC{}
\RightLabel{$\pi_2'''$}
\Deduce$!C, \Gamma \fCenter \Delta, !B$
\RightLabel{\CutCS}
\BinaryInf$\Gamma \fCenter \Delta, !B$
\DisplayProof
\]
La preuve $\pi_2'''$ est obtenue à partir de $\pi_2''$
par affaiblissement sans augmentation de hauteur.
Le rang de coupure est $<\cutr{\pi}$, puisque
$\depth{!C}<\depth{!B\lor!C}$.
L'hypothèse de récurrence fournit donc une !!{proof}
$\pi_1''$ dans \Log{G3c} de
$\Gamma\fCenter\Delta,!B$.
Considérons ensuite
\[
\AxiomC{}
\RightLabel{$\pi_1''$}
\Deduce$\Gamma \fCenter \Delta, !B$
\AxiomC{}
\RightLabel{$\pi_2'$}
\Deduce$!B, \Gamma \fCenter \Delta$
\RightLabel{\CutCS}
\BinaryInf$\Gamma \fCenter \Delta$
\DisplayProof
\]
Le rang de cette nouvelle coupure est
$=\depth{!B}<\depth{!B\lor!C}$.
La récurrence donne donc une !!{proof} $\pi'$
de $\Gamma\fCenter\Delta$.

\item $!A\ident(!B\land!C)$. Exercice.

\item Si $!A\ident(!B\lif!C)$, la preuve $\pi$ se termine par
\[
\AxiomC{}
\RightLabel{$\pi_1'$}
\Deduce$!B, \Gamma \fCenter \Delta, !C$
\RightLabel{\RightR{\lif}}
\UnaryInf$\Gamma \fCenter \Delta, !B \lif !C$
\LeftSubproofLabel{$\pi_1$}
\AxiomC{}
\RightLabel{$\pi_2'$}
\Deduce$\Gamma \fCenter \Delta, !B$
\AxiomC{}
\RightLabel{$\pi_2''$}
\Deduce$!C, \Gamma \fCenter \Delta$
\RightLabel{\LeftR{\lif}}
\BinaryInf$!B \lif !C, \Gamma \fCenter \Delta$
\RightSubproofLabel{$\pi_2$}
\RightLabel{\CutCS}
\BinaryInf$\Gamma \fCenter \Delta$
\DisplayProof
\]
La !!{proof} suivante se termine par une coupure de rang inférieur :
\[
\AxiomC{}
\RightLabel{$\pi_1'$}
\Deduce$!B, \Gamma \fCenter \Delta, !C$
\AxiomC{}
\RightLabel{$\pi_2''[!B \Sequent {}]$}
\Deduce$!C, !B, \Gamma \fCenter \Delta$
\RightLabel{\CutCS}
\BinaryInf$!B, \Gamma \fCenter \Delta$
\DisplayProof
\]
Dans la seconde prémisse, l'annotation
$\pi_2''[!B\Sequent{}]$ indique l'affaiblissement
de la preuve $\pi_2''$ par $!B$ à gauche, sans augmentation
de hauteur. Les deux prémisses ont ainsi le même contexte.
L'hypothèse de récurrence fournit une !!{proof}~$\pi_1''$
de $!B,\Gamma\fCenter\Delta$.
Considérons ensuite
\[
\AxiomC{}
\RightLabel{$\pi_2'$}
\Deduce$\Gamma \fCenter \Delta, !B$
\AxiomC{}
\RightLabel{$\pi_1''$}
\Deduce$!B, \Gamma \fCenter \Delta$
\RightLabel{\CutCS}
\BinaryInf$\Gamma \fCenter \Delta$
\DisplayProof
\]
Son rang de coupure est aussi plus petit que celui de~$\pi$.
La récurrence fournit donc une !!{proof}~$\pi'$
de $\Gamma\fCenter\Delta$ dans \Log{G3c}.

\item Si $!A\ident\lforall[x][!B(x)]$, la preuve $\pi$ se termine par
\[
\AxiomC{}
\RightLabel{$\pi_1'$}
\Deduce$\Gamma \fCenter \Delta, !B(c)$
\RightLabel{\RightR{\lforall}}
\UnaryInf$\Gamma \fCenter \Delta, \lforall[x][!B(x)]$
\LeftSubproofLabel{$\pi_1$}
\AxiomC{}
\RightLabel{$\pi_2'$}
\Deduce$!B(t), \lforall[x][!B(x)], \Gamma \fCenter \Delta$
\RightLabel{\LeftR{\lforall}}
\UnaryInf$\lforall[x][!B(x)], \Gamma \fCenter \Delta$
\RightSubproofLabel{$\pi_2$}
\RightLabel{\CutCS}
\BinaryInf$\Gamma \fCenter \Delta$
\DisplayProof
\]
L'hypothèse de récurrence permet de supprimer \CutCS de
\[
\AxiomC{}
\RightLabel{$\pi_1''$}
\Deduce$!B(t), \Gamma \fCenter \Delta, \lforall[x][!B(x)]$
\AxiomC{}
\RightLabel{$\pi_2'$}
\Deduce$!B(t), \lforall[x][!B(x)], \Gamma \fCenter \Delta$
\RightLabel{\CutCS}
\BinaryInf$!B(t), \Gamma \fCenter \Delta$
\DisplayProof
\]
La preuve $\pi_1''$ provient de $\pi_1$, et non de~$\pi_1'$,
par affaiblissement sans augmentation de hauteur.
Le rang de cette coupure est inchangé, mais sa hauteur
est plus petite. On obtient une !!{proof}~$\pi_2''$
de $!B(t),\Gamma\fCenter\Delta$.

La !!{proof}~$\pi_1'$ a pour séquent final
$\Gamma\Sequent\Delta,!B(c)$, où $c$ est la variable
propre d'une inférence~$\RightR{\lforall}$.
Ainsi, $c$ n'apparaît ni dans~$!B(x)$, ni dans
$\Gamma$, ni dans~$\Delta$.
Nous supposons la preuve~$\pi$ régulière :
la constante $c$ n'est la variable propre que de cette
inférence universelle, et elle apparaît seulement
au-dessus d'elle, donc dans~$\pi_1'$.
Elle n'est la variable propre d'aucune inférence
interne à cette sous-preuve.
Puisque $c$ n'apparaît pas dans~$\pi_2$, elle ne figure
pas dans~$t$. La régularité de toute la preuve garantit
aussi que ce terme de la seconde branche ne contient
aucune variable propre de la première.
Par le \olref[seq][qua]{lem:subst-regular},
$\Subst{\pi_1'}{t}{c}$ est une !!{proof} de
$\Gamma\Sequent\Delta,!B(t)$.
L'hypothèse de récurrence s'applique maintenant à
\[
\AxiomC{}
\RightLabel{$\Subst{\pi_1'}{t}{c}$}
\Deduce$\Gamma \fCenter \Delta, !B(t)$
\AxiomC{}
\RightLabel{$\pi_2''$}
\Deduce$!B(t), \Gamma \fCenter \Delta$
\RightLabel{\CutCS}
\BinaryInf$\Gamma \fCenter \Delta$
\DisplayProof
\]
Sa !!{formula} de coupure est~$!B(t)$ :
son rang est strictement inférieur à celui de~$\pi$.

\item Le cas $!A\ident\lexists[x][!B(x)]$ est laissé en exercice.
\end{enumerate}
\end{proof}

\begin{prob}
Vérifier le sous-cas du cas~D dans la preuve du
\olref{lem:cut-adm-G3c} où $!A$ n'est pas principal
dans la dernière inférence de~$\pi_1$, qui se termine
par~$\LeftR{\lif}$.
Une partie de l'exercice consiste à énoncer précisément
ce qu'il faut démontrer, c'est-à-dire la forme de $\pi$
dans ce cas.
\end{prob}

\begin{prob}
Vérifier le sous-cas du cas~C dans la preuve du
\olref{lem:cut-adm-G3c} où $!A$ n'est pas principal
dans la dernière inférence de~$\pi_1$, qui se termine
par~$\LeftR{\lforall}$.
\end{prob}

\begin{prob}
Vérifier le sous-cas du cas~E dans la preuve du
\olref{lem:cut-adm-G3c} où $!A$ est principal
dans les dernières inférences de $\pi_1$ et de $\pi_2$,
et où $!A\ident(!B\land!C)$.
\end{prob}

\begin{prob}
Vérifier le sous-cas du cas~E dans la preuve du
\olref{lem:cut-adm-G3c} où $!A$ est principal
dans les dernières inférences de $\pi_1$ et de $\pi_2$,
et où $!A\ident\lexists[x][!B(x)]$.
\end{prob}

Dans le cas~E, la hauteur de coupure à laquelle on
applique la récurrence peut dépasser celle de la preuve
initiale. Par exemple, pour $!A\ident(!B\lif!C)$,
nous avons transformé $\pi$ en une !!{proof} avec
deux inférences \CutCS :
\[
\AxiomC{}
\RightLabel{$\pi_2'$}
\Deduce$\Gamma \fCenter \Delta, !B$
\AxiomC{}
\RightLabel{$\pi_1'$}
\Deduce$!B, \Gamma \fCenter \Delta, !C$
\AxiomC{}
\RightLabel{$\pi_2''[!B \Sequent {}]$}
\Deduce$!C, !B, \Gamma \fCenter \Delta$
\RightLabel{\CutCS}
\BinaryInf$!B, \Gamma \fCenter \Delta$
\RightSubproofLabel{$\pi_1''$}
\RightLabel{\CutCS}
\BinaryInf$\Gamma \fCenter \Delta$
\DisplayProof
\]
La coupure supérieure, entre $\pi_1'$ et la preuve
affaiblie $\pi_2''$, a un rang et une hauteur de coupure
plus petits que ceux de~$\pi$.
Mais la !!{proof}~$\pi_1''$ obtenue par récurrence
peut avoir une hauteur supérieure.
La coupure inférieure porte toutefois sur la
!!{formula} $!B$, et son \emph{rang} reste inférieur
à celui de~$\pi$. C'est pourquoi la récurrence
lexicographique privilégie le rang.

Le \olref{lem:cut-adm-G3c} permet de supprimer une unique
inférence \CutCS la plus haute d'une !!{proof}.
On peut éliminer un nombre quelconque de telles
inférences en appliquant successivement le lemme
aux sous-!!{proof}s qui se terminent par une coupure
la plus haute.

\begin{thm}\ollabel{thm:cut-elim-G3c-topmost}
Toute !!{proof} $\pi$ dans $\Log{G3c}+\CutCS$
peut être transformée en une !!{proof} dans~\Log{G3c}
du même séquent final.
\end{thm}

\begin{proof}
Raisonnons par récurrence sur le nombre~$n$
d'inférences \CutCS dans~$\pi$.
Si $n=0$, la preuve est déjà dans~\Log{G3c}.
Sinon, choisissons la sous-!!{proof}~$\pi'$
qui se termine par une coupure la plus haute.
Elle ne contient qu'une seule inférence \CutCS,
qui est sa dernière inférence.
Par le \olref{lem:cut-adm-G3c}, on la transforme
en une preuve~$\pi''$ sans coupure.
Remplaçons $\pi'$ par $\pi''$ dans la preuve entière :
la nouvelle !!{proof} contient $n-1$ inférences \CutCS,
et l'hypothèse de récurrence s'applique.
\end{proof}

\end{document}
